\documentclass[11pt]{amsart}
\usepackage{amsmath,amssymb,amsthm,mathtools}
\usepackage{enumitem}

%%%%%%%%%%%%%%%%%%%%%%%%%%%%%%%%%%%%%%%%%%%%%%%%%%%%%%%%%%%%
\title{Transport Theory of Observation\\
A Mathematical Framework for Receipts, Fibers, Kernels and Reconstruction}
\author{Jean-Marie Tassy}
%%%%%%%%%%%%%%%%%%%%%%%%%%%%%%%%%%%%%%%%%%%%%%%%%%%%%%%%%%%%

\newtheorem{definition}{Definition}[section]
\newtheorem{theorem}[definition]{Theorem}
\newtheorem{lemma}[definition]{Lemma}
\newtheorem{proposition}[definition]{Proposition}
\newtheorem{corollary}[definition]{Corollary}
\newtheorem{remark}[definition]{Remark}
\newtheorem{example}[definition]{Example}

%%%%%%%%%%%%%%%%%%%%%%%%%%%%%%%%%%%%%%%%%%%%%%%%%%%%%%%%%%%%
\begin{document}
\maketitle

\begin{abstract}
We introduce a general mathematical theory of observation.
Rather than studying a system directly, we study the map from reality to
observable information.
This produces natural notions of
observation maps,
receipt fibers,
generalized kernels,
observable quotient spaces,
transport operators,
and reconstruction.
The framework applies independently of any specific application and naturally
connects to topology, inverse problems, information theory,
control theory, statistics, and quantum measurement.
\end{abstract}

%%%%%%%%%%%%%%%%%%%%%%%%%%%%%%%%%%%%%%%%%%%%%%%%%%%%%%%%%%%%
\section{Observation Maps}

Let $\mathcal{S}$ be a state space.
Let $\mathcal{L}$ be an observation space.

\begin{definition}[Observation Map]
An observation map is a function
\[
R: \mathcal{S} \rightarrow \mathcal{L}.
\]
The value $R(s)$ is called the receipt (or observation) associated with
the state $s$.
\end{definition}

The map need not be
\begin{itemize}
\item linear,
\item continuous,
\item differentiable,
\item probabilistic.
\end{itemize}
It is purely an observation operator.

%%%%%%%%%%%%%%%%%%%%%%%%%%%%%%%%%%%%%%%%%%%%%%%%%%%%%%%%%%%%
\section{Receipt Fibers}

Observation identifies states that appear identical.

\begin{definition}
The fiber over $\ell \in \mathcal{L}$ is
\[
R^{-1}(\ell)
=
\{ s \in \mathcal{S} : R(s) = \ell \}.
\]
\end{definition}

Equivalently, the fiber through a state $s$ is
\[
[s]_R
=
\{ x \in \mathcal{S} : R(x) = R(s) \}.
\]

%%%%%%%%%%%%%%%%%%%%%%%%%%%%%%%%%%%%%%%%%%%%%%%%%%%%%%%%%%%%
\section{Observable Equivalence}

\begin{definition}
Define $s_1 \sim_R s_2$ iff $R(s_1) = R(s_2)$.
\end{definition}

\begin{lemma}
The relation $\sim_R$ is an equivalence relation.
\end{lemma}

\begin{proof}
Immediate from equality.
\end{proof}

%%%%%%%%%%%%%%%%%%%%%%%%%%%%%%%%%%%%%%%%%%%%%%%%%%%%%%%%%%%%
\section{Observable Quotient Space}

\begin{definition}
The observable quotient is
\[
\mathcal{S} / {\sim_R}.
\]
Its elements are the receipt fibers.
\end{definition}

This quotient is the entire observable universe induced by the observer.

%%%%%%%%%%%%%%%%%%%%%%%%%%%%%%%%%%%%%%%%%%%%%%%%%%%%%%%%%%%%
\section{Generalized Kernel}

Classical kernels belong to linear algebra.
Observation suggests a broader concept.

\begin{definition}
A transformation $T: \mathcal{S} \rightarrow \mathcal{S}$ is invisible to
the observation map whenever
\[
R \circ T = R.
\]
The collection of all such transformations is
\[
\operatorname{Inv}(R)
=
\{ T : R \circ T = R \}.
\]
\end{definition}

We call $\operatorname{Inv}(R)$ the generalized kernel (or invariance monoid)
of the observation.

%%%%%%%%%%%%%%%%%%%%%%%%%%%%%%%%%%%%%%%%%%%%%%%%%%%%%%%%%%%%
\section{Invisible Transformations}

\begin{definition}
A transformation is observable if $R(T(s)) \neq R(s)$ for some state.
Otherwise it is invisible.
\end{definition}

Invisible transformations move inside fibers.
Observable transformations move between fibers.

%%%%%%%%%%%%%%%%%%%%%%%%%%%%%%%%%%%%%%%%%%%%%%%%%%%%%%%%%%%%
\section{Reconstruction}

Observation naturally produces an inverse problem.

\begin{definition}
Given $\ell \in \mathcal{L}$, the reconstruction problem is to determine
$R^{-1}(\ell)$.
\end{definition}

%%%%%%%%%%%%%%%%%%%%%%%%%%%%%%%%%%%%%%%%%%%%%%%%%%%%%%%%%%%%
\section{Faithfulness}

\begin{definition}
The observation map is faithful if
\[
R(s_1) = R(s_2) \Longrightarrow s_1 = s_2.
\]
Equivalently, $R$ is injective.
\end{definition}

%%%%%%%%%%%%%%%%%%%%%%%%%%%%%%%%%%%%%%%%%%%%%%%%%%%%%%%%%%%%
\section{Transport Theorem}

\begin{theorem}[Non-Reconstructibility]
If $R$ is not injective, then perfect reconstruction from observations
is impossible.
\end{theorem}

\begin{proof}
If $R(s_1) = R(s_2)$, $s_1 \neq s_2$, then one observation corresponds
to at least two states. No reconstruction algorithm can uniquely determine
which state generated the observation.
\end{proof}

%%%%%%%%%%%%%%%%%%%%%%%%%%%%%%%%%%%%%%%%%%%%%%%%%%%%%%%%%%%%
\section{Fiber Geometry}

The geometry of fibers determines the information lost by observation.
Large fibers imply more ambiguity.
Singleton fibers imply perfect identification.

\begin{definition}[Fiber Diameter]
Given a secondary metric $d$ on $\mathcal{S}$, the fiber diameter is
\[
\operatorname{diam}([s]_R, d)
=
\sup \{ d(x, y) : x, y \in [s]_R \}.
\]
A complete observation map has $\operatorname{diam}([s]_R, d) = 0$ for
all $s$. A non-faithful map has strictly positive diameter in some fiber.
\end{definition}

\begin{remark}
The diameter is not observable from $R$ alone. It requires a secondary
observation map $P: \mathcal{S} \rightarrow \mathcal{M}$ to induce a metric
$d_P(x, y) = d_{\mathcal{M}}(P(x), P(y))$. The fiber geometry under
$d_P$ measures what $P$ sees that $R$ misses.
\end{remark}

%%%%%%%%%%%%%%%%%%%%%%%%%%%%%%%%%%%%%%%%%%%%%%%%%%%%%%%%%%%%
\section{Transport Operators}

Given two observation maps $R_1: \mathcal{S} \to \mathcal{L}_1$ and
$R_2: \mathcal{S} \to \mathcal{L}_2$, one may ask whether there exists
\[
T: \mathcal{L}_1 \rightarrow \mathcal{L}_2
\]
such that $T \circ R_1 = R_2$.

This defines transport between observation systems.

\begin{definition}[Inter-Fiber Transport]
Given two fibers $[s]_R$ and $[s']_R$ in the same state space, a transport
operator is a map
\[
\tau: [s]_R \longrightarrow [s']_R
\]
that moves between fibers while preserving some specified structure.
\end{definition}

This connects to parallel transport, gauge transport, optimal transport
between probability distributions, and functorial semantics.

%%%%%%%%%%%%%%%%%%%%%%%%%%%%%%%%%%%%%%%%%%%%%%%%%%%%%%%%%%%%
\section{Completeness}

\begin{definition}
An observation system is complete whenever every state is uniquely determined
by its receipt. Equivalently, $R$ is injective.
\end{definition}

\begin{definition}[Sufficient Observation]
$R$ is sufficient for a parameter $\theta: \mathcal{S} \to \Theta$ if
\[
s_1 \sim_R s_2 \Longrightarrow \theta(s_1) = \theta(s_2).
\]
That is, the receipt determines the parameter, even if it does not
determine the state.
\end{definition}

This connects to the Lehmann--Scheff\'{e} theorem in statistics:
a sufficient statistic need not be complete (injective) to uniquely
determine a parameter.

%%%%%%%%%%%%%%%%%%%%%%%%%%%%%%%%%%%%%%%%%%%%%%%%%%%%%%%%%%%%
\section{Examples}

The abstraction appears naturally in many domains.

\begin{example}[Governance receipts]
$\mathcal{S}$ = system states. $R$ = ledger receipt chain.
Faithful iff the receipt chain uniquely determines the system state.
The E11/E12 parallel-session fork is an explicit instance of
$R(s_1) = R(s_2)$ with $s_1 \neq s_2$.
\end{example}

\begin{example}[Authority language linter]
$\mathcal{S}$ = documents. $R(s)$ = warning set output by the linter.
The fiber $[s]_R$ over the empty warning set contains all documents the
linter cannot distinguish from clean text. The ``semantic gap'' is the
fiber diameter under a secondary semantic embedding.
\end{example}

\begin{example}[Riemann Hypothesis]
$\mathcal{S}$ = zero configurations $Z$. $R_J(Z)$ = finite-band positivity
operator applied to $Z$. Faithfulness question: does $R_J(Z_1) = R_J(Z_2)$
with $Z_1 \neq Z_2$? If yes, finite-band positivity cannot characterize RH.
\end{example}

\begin{example}[Embedding models]
$\mathcal{S}$ = semantic meanings. $E: \mathcal{S} \to \mathbb{R}^d$.
$\operatorname{Inv}(E)$ contains paraphrase, syntax variation, and negation
in context. Characterizing $\operatorname{Inv}(E)$ is an active research problem.
\end{example}

\begin{example}[Control theory]
$\mathcal{S}$ = system states. $R$ = sensor readings. The unobservable
subspace is the fiber. Observability = faithfulness of $R$.
\end{example}

\begin{example}[Quantum measurement]
$\mathcal{S}$ = quantum states. $R$ = measurement operator.
States modulo global phase occupy the same fiber.
The observable universe is the projective space $\mathcal{S} / {\sim_R}$.
\end{example}

%%%%%%%%%%%%%%%%%%%%%%%%%%%%%%%%%%%%%%%%%%%%%%%%%%%%%%%%%%%%
\section{Research Directions}

The framework immediately suggests several mathematical programs.

\begin{enumerate}[label=\arabic*.]
\item Classification of observation maps by fiber structure.
\item Geometry and topology of receipt fibers.
\item Algebra of invisible transformations ($\operatorname{Inv}(R)$ as group/monoid).
\item Information loss invariants: entropy of the quotient map.
\item Minimal complete receipt systems (analog of minimal sufficient statistics).
\item Transport between observation categories (functorial).
\item Entropy of observation: $H(S) - H(S \mid R(S))$.
\item Stability of fibers under perturbation of $R$.
\item Reconstruction complexity: cost of moving within a fiber.
\item Fiber metric axioms: which secondary maps $P$ induce geometrically meaningful fiber distances?
\end{enumerate}

%%%%%%%%%%%%%%%%%%%%%%%%%%%%%%%%%%%%%%%%%%%%%%%%%%%%%%%%%%%%
\section{Central Principle}

Observation does not reveal reality.
Observation reveals an equivalence class.

\[
\boxed{
\text{Reality}
\longrightarrow
\text{Observation}
\longrightarrow
\text{Fiber}
}
\]

The observer never directly accesses $\mathcal{S}$, but only the quotient
$\mathcal{S} / {\sim_R}$.

%%%%%%%%%%%%%%%%%%%%%%%%%%%%%%%%%%%%%%%%%%%%%%%%%%%%%%%%%%%%
\end{document}
